% No preamble. Prefix fcc-. Actual finite references and finite periodic backgrounds.
\section{Finite-regulator classical contraction}
\label{fcc-section}

Fix dimension four, a compact simple gauge group with the fixed
color normalization, a primitive odd ratio \(\ell\), and
\(0<\alpha<1\). A semisimple product uses one reference coefficient
per simple factor. Fix the activity exponent \(a_0>0\) and a
grouping \(J\). Let \(L\) be an allowed endpoint ratio and write
\begin{equation}
 P=LN.
 \label{fcc-periods}
\end{equation}
Here \(P\) is the input current-unit torus period and \(N\) is
the output current-unit period. Both cubic tori are compatible
with the partitions; in particular \(J\mid N\). The preceding
homogeneous depth \(M\) is arbitrary.

On a current torus of period \(Q\), let
\begin{equation}
 \|E\|_{Q;a_0,b}
   =\sup_i\sum_{H\ni i}e^{a_0|H|}\|E_H\|_b,
 \qquad
 \Phi_E(A)=\sum_H\bigl[E_H(A|_H)-E_H(0)\bigr],
 \label{fcc-norm-total}
\end{equation}
where \(H\) ranges over nonempty connected sets of \(J\)-cells
and \(\|\cdot\|_b\) is the supremum norm on the complex
scalar-coordinate \(b\)-polydisc. Quotient by equality of
\(\Phi_E\) on that full polydisc, and denote the quotient norm by
\(\|[E]\|_Q\). The equality kernel is closed, since finite-volume
evaluations are continuous, so this is a Banach quotient.
We use its invariant subspace: local gauge, translation,
hypercubic and global color identities are imposed on totals.
Individual coefficients need not be invariant. Vacuum constants
are tracked separately.

\subsection{Actual finite inputs and the reference projection}
\label{fcc-inputs}

The earlier finite constructions supply the following inputs.
First, for \(n=2,3,4\), fixed-order Cauchy estimates and connected
carriers give, for some fixed \(\mu=\mu(J)>0\),
\begin{equation}
 \sup_{b_1}\sum_{b_2,\ldots,b_n}
 e^{\mu D_Q(b_1,\ldots,b_n)}
 |K_{n,Q}(b_1,\ldots,b_n)|
 \le C_n(J)b^{-n}\|E\|_{Q;a_0,b},
 \qquad K_{n,Q}=D^n\Phi_E(0).
 \label{fcc-kernel-bound}
\end{equation}
Here \(D_Q\) is the maximum pairwise torus \(\ell^1\) distance
between bond midpoints, and internal indices are included in
the sums. Indeed, every differentiated tuple lies in a common
carrier, \(D_Q\le C J|H|\), while its fixed-degree coordinate
count is a polynomial in \(D(J)|H|\). Choosing \(\mu\) small
and spending the remaining cardinal reserve absorbs this
polynomial without a volume factor.

Second, \(S_{M,Q}\) is the actual finite homogeneous classical
reference. It has zero value and first derivative at zero and
the same total symmetries. Its strict reference family on one
fixed radius \(\rho\) is uniformly bounded by \(M_S\). Its flat
Hessian is the exact relative periodization of the uniformly
exponentially decaying \(q_M\), whose midpoint Maxwell
coefficient is one.

Removing coefficientwise homogeneous degrees zero and one
does not change its normalized total. The phase projection
bound gives, for \(0<b\le\rho/2\),
\begin{equation}
 \|S_{M,Q}^{(\ge2)}\|_{Q;a_0,b}
 \le \frac{M_S(b/\rho)^2}{1-b/\rho}
 \le C_S b^2,\qquad C_S=2M_S/\rho^2.
 \label{fcc-reference-small}
\end{equation}
This does not assert that the original individual reference
coefficients have zero linear jets.

For \(Q\ge64\), symmetrically lift the quadratic torus kernel
using midpoint diameter at most \(\lfloor Q/16\rfloor\), and
write its midpoint symbol as \(\Gamma^E_{2,Q}\). For distinct
directions \(\mu,\nu\), with the invariant color normalization,
put
\begin{equation}
 b_Q(E)=\frac12
       \partial_{p_\nu}^2\Gamma^E_{2,Q;\mu\mu}(0).
 \label{fcc-cutoff-jet}
\end{equation}
Point symmetry makes the distinct pair immaterial. The exact
flat reference periodization implies
\begin{equation}
 |b_Q(S_{M,Q})-1|\le C(J)e^{-c(J)Q}.
 \label{fcc-reference-denominator}
\end{equation}
This estimate includes both the discarded infinite kernel
tail and its long images folded into the short torus lift.
Uniformity of the \(q_M\) decay makes it uniform in \(M\).
Choose \(Q_0(J)\) so that the denominator has modulus at least
\(1/2\), and define
\begin{equation}
 \beta_{M,Q}([E])
     =\frac{b_Q(E)}{b_Q(S_{M,Q})},
 \qquad
 \mathcal P_{M,Q}[E]
     =[E]-\beta_{M,Q}([E])[S_{M,Q}].
 \label{fcc-projection}
\end{equation}
The kernels depend only on the normalized total, so these
are well defined on the quotient. Equations
\eqref{fcc-kernel-bound}--\eqref{fcc-reference-small} imply
\begin{align}
 |\beta_{M,Q}([E])|
     &\le C_\beta(J)b^{-2}\|[E]\|_Q,\notag\\
 \beta_{M,Q}([S_{M,Q}])&=1,\qquad
 \mathcal P_{M,Q}^2=\mathcal P_{M,Q},\qquad
 \|\mathcal P_{M,Q}\|\le C_{\rm ext}(J)
       :=1+C_\beta(J)C_S.
 \label{fcc-projection-bound}
\end{align}
The projection bound, unlike the bound on \(\beta\) alone,
is independent of the small \(b\).

Third, the actual finite pre-gauge ordinary representative
\(a=a_{M,L,N}(w)\) is the partial ordinary average of the
full \(ML\)-refined ordinary minimizer. The exact finite
two-minimum comparison gives
\begin{equation}
 a(0)=0,\qquad
 S_{M,LN}(a(w))=S_{ML,N}(w).
 \label{fcc-reference-transport}
\end{equation}
The identity first holds for real small fields and then
holomorphically. In dimension four it has no additional
dimensional prefactor.

For any fixed finite cardinal reserve \(A\), this map has
coherent complete-field support coefficients on final unit
cells:
\begin{align}
 a|_i&=\sum_{\varnothing\ne X}a_{i,X}(w_X),
       \qquad a_{i,X}(0)=0,\notag\\
 \sup_i\sum_X e^{\kappa\tau_N(\{i\}\cup X)+A|X|}
       \|a_{i,X}\|_{C^{1,\alpha};b}
      &\le C_* b/L.
 \label{fcc-pregauge-family}
\end{align}
The whole current-field output fiber and its physical
gradient H\"older quotients are inside the support sum.
Coefficients on overlapping fixed collars are restrictions
of one global assignment. The constants and the permitted
classical radius are uniform in \(M,L,N\). This is the
finite complete-field prefix theorem, not a periodization
assumption about an infinite background.

Define scalar classical substitution by
\begin{equation}
 (\mathcal T_{M,L,N}\Phi_E)(w)=\Phi_E(a_{M,L,N}(w)).
 \label{fcc-transport}
\end{equation}
For \(L\) sufficiently large, the amplitude is strictly
inside the input \(b\)-domain, so equality of finite totals
implies equality after substitution. The estimates below
give a bounded output representative.

We also record the finite-domain invariance argument.
For a real output field and a sufficiently small local
gauge transform, the finite transfer and gauge-flow
construction gives an actual near-identity connecting
input gauge \(g\). Both endpoint angles have size
\(\rho_0=C_A b/L\). The bond equation \(V^g=U\) bounds
the edge increments of \(g\) by
\(C(|U-I|+|V-I|)\), and hence
\(|D\log g|\le C\rho_0\) on its fixed near-identity chart.
The existence of the small chart comes from that actual
finite construction, not merely from an edge bound.
Every field along \(\exp(t\log g)\), \(0\le t\le1\), has
scalar-coordinate size at most \(C_g\rho_0\). Fixed
matrix/color norm conversions are included in \(C_g\).
Taking \(L>C_gC_A\) with a strict margin keeps the entire
path inside \(b\). The input Ward identity integrates
along it, and analytic continuation gives the local
complex output Ward identity. No infinite gauge
truncation is needed. The finite minimum comparisons
or exact covariance likewise give the spatial and
global color identities on this common domain.

Consequently \(\mathcal T\) acts on the finite invariant
quotients. The full-reference identity gives exactly
\begin{equation}
 \mathcal T_{M,L,N}[S_{M,LN}]=[S_{ML,N}],\qquad
 \mathcal P_{ML,N}\mathcal T_{M,L,N}\mathcal P_{M,LN}
       =\mathcal P_{ML,N}\mathcal T_{M,L,N}.
 \label{fcc-telescope}
\end{equation}
This does not assert that \(\mathcal T\) preserves the
zero-\(\beta\) subspace.

\subsection{Graded Ward defects and finite-range jet repair}
\label{fcc-repair-section}

Let \([F]\) satisfy \(\beta_{M,P}([F])=0\), and choose a
representative of norm at most \(\varepsilon\), as close
to its quotient norm as desired. Set \(T_0=\lfloor P/16\rfloor\).
The symmetric infinite lift
\begin{equation}
 \widetilde K_n(b_1,\ldots,b_n)
  =K_{n,P}(b_1\bmod P,\ldots,b_n\bmod P)
           {\bf1}_{\{D_\infty(b_1,\ldots,b_n)\le T_0\}}
 \label{fcc-short-lift}
\end{equation}
preserves slot symmetry, midpoint parity and the spatial
and color symmetries. A short torus tuple has a unique
lift with its first bond fixed: all coordinate differences
are below half a period, and the lifted pairwise distances
are the torus distances. Thus
\begin{equation}
 K_{n,P}=\operatorname{Per}_P\widetilde K_n+K_n^{\rm tail},
 \qquad
 K_n^{\rm tail}=K_{n,P}{\bf1}_{\{D_P>T_0\}}.
 \label{fcc-short-tail}
\end{equation}
Every fixed spatial moment of the tail is bounded by
\(C(J,m)\varepsilon b^{-n}e^{-c(J)P}\).

Ward operations are finite sums of nearest-neighbor
shifts, coincident-slot insertions and color commutators.
They commute with relative periodization. A lifted Ward
defect has diameter at most \(T_0+2\), where periodization
is injective. Its infinite fixed moments are therefore
bounded by those of the discarded tails.

Keep the field-degree grading. For \(b\le1\), the
degree-\(n\) identity involves only degrees \(2,\ldots,n\),
so its fixed-moment error is bounded by
\begin{equation}
 \delta_n=C(J)\varepsilon b^{-n}e^{-c(J)P},
       \qquad n=2,3,4.
 \label{fcc-graded-defect}
\end{equation}
The quadratic coefficient elimination uses only
\(\delta_2\); the cubic uses \(\delta_3\) and the already
bounded lower jets, and the quartic uses \(\delta_4\)
and its lower jets. Since \(b\le1\), these lower terms
fit the displayed grading. A common \(b^{-4}\) estimate
for all degrees is neither needed nor used.

For completeness, the lifted hypercubic Hessian has
mass \(m\) and quadratic coefficients \(a,b,c\).
The degree-one and degree-three Ward comparisons bound
\(m,a,c\) by \(C\delta_2\), because the extracted
coefficient \(b_P(F)\) is exactly zero. Odd midpoint
orders vanish by parity. Write the cubic linear polynomial
as \(A p+B q\); the independent \(p_\mu q_\sigma\) Ward
coefficients bound \(B\), and the \(q_\nu p_\sigma\)
coefficients bound \(A\), by \(C\delta_3\). The quartic
leading identity contracts its constant tensor with an
arbitrary gauge momentum and bounds it by
\(C\delta_4\). Hence the lifted quadratic jets through
degree three, cubic jets through degree one, and quartic
constant are respectively \(O(\delta_2),O(\delta_3),
O(\delta_4)\).

Pass to tail-coordinate Laurent symbols \(F_n(z)\),
anchoring the first positive-bond tail and retaining
internal indices. Fixed midpoint phases and the local
change \(z=e^{ip}\) preserve these finite-order bounds.
Put \(m_2=3,m_3=1,m_4=0\), and take
\begin{equation}
 J_{m_n}F_n(z)=
    \sum_{|\gamma|\le m_n}
       \frac{\partial_z^\gamma F_n(1)}{\gamma!}
                         (z-1)^\gamma.
 \label{fcc-laurent-jet}
\end{equation}
Its kernel is of fixed finite relative range and has
absolute row norm at most \(C\delta_n\).

Average this kernel over the \(n!\) slot permutations,
reanchoring the first slot each time, and call the result
\(R_n\). Its relative range increases only by a fixed
factor and its norm remains \(O(\delta_n)\).
The corrected kernel
\begin{equation}
 \widehat K_n=\widetilde K_n-R_n
 \label{fcc-repaired-kernel}
\end{equation}
has exactly zero jets through degree \(m_n\). Indeed
each reanchored residual is the composition of
\(F_n-J_{m_n}F_n\) with an analytic monomial change
fixing \(z=1\), possibly times a harmless analytic
phase. This preserves vanishing order \(m_n+1\);
averaging preserves it as well. Reality and color
covariance can be retained. Point symmetry of the
repair is not required by the spatial estimate.

For a fixed smaller \(\sigma=\sigma(J)>0\),
\begin{align}
 \sup_{\text{first index}}\sum_{\text{others},\,r_2,\ldots,r_n}
 e^{\sigma\sum_{j\ge2}|r_j|}
      |\widehat K_n(0,r_2,\ldots,r_n)|
       &\le C(J)\varepsilon b^{-n},\notag\\
 \|R_n\|_{\rm row}
       &\le C(J)\varepsilon b^{-n}e^{-c(J)P}.
 \label{fcc-repaired-bounds}
\end{align}
The second line has fixed finite relative range.
The finite polynomial decomposition is exact:
\begin{equation}
 \frac{K_{n,P}(a,\ldots,a)}{n!}
  =\frac{(\operatorname{Per}_P\widehat K_n)(a,\ldots,a)}{n!}
   +\frac{(\operatorname{Per}_P R_n)(a,\ldots,a)}{n!}
   +\frac{K_n^{\rm tail}(a,\ldots,a)}{n!}.
 \label{fcc-polynomial-decomposition}
\end{equation}
No error term is discarded or separately asserted to
be gauge invariant.

\subsection{Direct periodic spatial localization}
\label{fcc-periodic-spatial-section}

We give the finite version of the spatial argument explicitly.
For a relative infinite kernel \(\widehat K_n\), its
periodized polynomial equals
\begin{equation}
 \frac1{n!}\sum_{x\in T_P}
  \sum_{r_2,\ldots,r_n\in\mathbb Z^4}
   \widehat K_n(r_2,\ldots,r_n)
       [a(x),a(x+r_2),\ldots,a(x+r_n)],
 \label{fcc-periodic-polynomial}
\end{equation}
where fields are read periodically. The relative sums
converge absolutely, and repeated torus slots retain
their multilinear multiplicity.

For the quadratic symbol \(F(z)=\sum_r\widehat K_2(r)z^r\),
define
\[
 B_jF(z)=
 \frac{F(1,\ldots,1,z_j,\ldots,z_d)
       -F(1,\ldots,1,1,z_{j+1},\ldots,z_d)}{z_j-1}.
\]
The quotient is removable, and
\(F-F(1)=\sum_j(z_j-1)B_jF\). A monomial contributes
\(|r_j|\) terms of exponent length at most \(|r|_1\).
Thus the combined row Wiener norm, including the \(j\)
sum, loses at most \([e(\sigma-\sigma')]^{-1}\).
Vanishing order decreases by at most one. Apply twice
and put \(\widetilde H_{\rho\tau}=-z_\rho B_\tau B_\rho F\).
Then
\begin{equation}
 F=\sum_{\rho,\tau}(z_\rho^{-1}-1)
          \widetilde H_{\rho\tau}(z)(z_\tau-1),
 \qquad \widetilde H_{\rho\tau}(1)=0.
 \label{fcc-divergence-factor}
\end{equation}
The combined row norm at exponent \(\sigma/2\) is at
most
\[
 \frac{16e^{\sigma/2}}{e^2\sigma^2}
                  C(J)\varepsilon b^{-2}.
\]

Finite periodic summation by parts has no boundary term.
Equation \eqref{fcc-divergence-factor} and the zero sum
of \(\widetilde H\) give exactly
\begin{equation}
 \frac12\sum_{\substack{x\in T_P,\,r\\\rho,\tau}}
 [D_\rho a(x)]^T\widetilde H_{\rho\tau}(r)
                [D_\tau a(x+r)-D_\tau a(x)].
 \label{fcc-periodic-ibp}
\end{equation}
No summability of an infinite field is assumed.

For the cubic term set
\(\ell_r(x)=\sum_\rho r_\rho D_\rho a(x)\) and
\(\operatorname{Rem}_r=a(x+r)-a(x)-\ell_r(x)\).
The exact lifted constant and first moments cancel
the all-anchor and single-\(\ell\) terms. The six
remaining terms are the two \(a,a,\operatorname{Rem}\)
terms, \(a,\ell,\ell\), the two
\(a,\ell,\operatorname{Rem}\) terms, and
\(a,\operatorname{Rem},\operatorname{Rem}\).
For the quartic term subtract the all-anchor product
and telescope its three shifted factors. These are
identities for periodic fields with lifted integer
offsets; no periodic affine function is introduced.

Let \(h=L^{-1}\) and \(m=\sum_{j\ge2}|r_j|_1\).
For \(m\le L\), coherence in
\eqref{fcc-pregauge-family} gives, in the weighted
family norm on one fixed enlarged collar,
\begin{align}
 |a|&\le Cbh,& |Da|&\le Cbh^2,\notag\\
 |Da(x+r)-Da(x)|&\le Cbh^{2+\alpha}|r|^\alpha,&
 |\operatorname{Rem}_r|&\le
                 Cbh^{2+\alpha}|r|^{1+\alpha}.
 \label{fcc-near-estimates}
\end{align}
Negative path traversals use the positive difference
at the lower endpoint. Products of coefficient families
glue their support trees and preserve cardinal reserve.
The quadratic, cubic and quartic residual products
are therefore bounded respectively by
\[
 Cb^2h^{4+\alpha}|r|^\alpha,\qquad
 Cb^3h^{4+\alpha}(1+m)^6,\qquad
 Cb^4h^5(1+m).
\]
The kernel moments cost only fixed powers of
\(\sigma^{-1}\), and their \(b^{-n}\) cancels \(b^n\).

For every product, project the finite endpoint root
set and every terminal set \(\{\text{root}\}\cup X\)
to final \(J\)-cells. Choose the short connecting
graphs supplied by the projection lemma and take
their connected union \(H\). The carrier need not
contain the full projection of an arbitrarily chosen
fine tree. With a fixed number \(n_0\) of endpoint
roots,
\begin{align}
 |H|&\le1+T_H,\notag\\
 T_H&\le\frac2J\left[
          \tau_N(\text{roots})
           +\sum_{\rm occ}\tau_N(\{\text{root}\}\cup X)\right]
             +4\left[n_0+\sum_{\rm occ}|X|\right].
 \label{fcc-carrier-geometry}
\end{align}
Here \(\tau_N\) uses final unit cells of period \(N\).
Shortest periodic paths obey the same bound; wrapping
can only reduce the union cardinality. If \(j(x)\)
is the final \(J\)-cell containing \(x/L\), then
\begin{equation}
 e^{a_0|H|}{\bf1}_{\{z\in H\}}
 \le e^{a_0}e^{-\operatorname{dist}_N(z,j(x))}
                         e^{(a_0+1)T_H}.
 \label{fcc-reroot}
\end{equation}
The distance here is on final \(J\)-cells, of period
\(N/J\), not on unit cells. Choose
\begin{equation}
 A\ge4(a_0+1),\qquad 2(a_0+1)/J\le\kappa.
 \label{fcc-reserve}
\end{equation}
These pay the support and attached-tree weights.
Finally,
\begin{equation}
 \sum_{x\in T_P}e^{-\operatorname{dist}_N(z,j(x))}
       \le(JL)^4 S_1,\qquad
 S_1=\sum_{y\in\mathbb Z^4}e^{-|y|_1}.
 \label{fcc-anchor-count}
\end{equation}
There are exactly \((JL)^4\) current anchor sites per
output \(J\)-cell. Choosing a shortest representative
per torus coset bounds the remaining sum by \(S_1\).

For \(m>L\), use separate endpoint roots. Their
joining tree has length at most \(C(m/L+1)\), so
\eqref{fcc-reroot} adds at most
\(C e^{C'm/(JL)}\). The rough quadratic, cubic and
quartic residual bounds are respectively
\(Cb^2h^4\), \(Cb^3h^3(1+m)^2\), and \(Cb^4h^4\).
Taking \(L\) large after \(J\) makes \(C'/(JL)\le\sigma/8\).
The unspent kernel decay bounds all far terms by
\[
 C(J)\varepsilon(L+2)e^{-\sigma L/8}
       \le C'(J)\varepsilon L^{-\alpha}.
\]
The near terms have the same or better bound after
\eqref{fcc-anchor-count}. Normal regrouping by \(H\)
therefore constructs a connected output family with
\begin{equation}
 \left\|\sum_{n=2}^4
       \frac{(\operatorname{Per}_P\widehat K_n)
                       (a,\ldots,a)}{n!}\right\|_{\rm rep}
       \le C_{\rm low}(J)\varepsilon L^{-\alpha}.
 \label{fcc-canceled-estimate}
\end{equation}
The notation \(\|\cdot\|_{\rm rep}\) denotes the norm
of the constructed representative, hence bounds its
quotient norm. Its coefficients are holomorphic on
the same \(b\)-polydisc. This proof used the actual
finite periodic support family, not a putative
infinite-volume extension of it.

\subsection{Wrapping errors at the unchanged activity exponent}
\label{fcc-errors-section}

We use a direct endpoint-product bound. Suppose a torus
\(n\)-kernel \(W\), \(n\le4\), satisfies
\begin{equation}
 \sup_{b_1}\sum_{b_2,\ldots,b_n}e^{\mu D_P}|W|
                         \le A_n.
 \label{fcc-error-kernel}
\end{equation}
Join each endpoint to the first by a shortest torus
path. Its root tree on final unit cells has length
at most \(C_n(D_P/L+1)\). Insert the actual endpoint
support families and use
\eqref{fcc-carrier-geometry}--\eqref{fcc-reroot}.
The additional factor is at most
\[
 C\exp\!\left(\frac{C_n'(a_0+1)D_P}{JL}\right).
\]
Increase the lower threshold for \(L\) until this
exponent coefficient is at most \(\mu/4\).
The field product contributes \((Cb/L)^n\), while
\eqref{fcc-anchor-count} gives
\begin{align}
 \left\|\frac{W(a,\ldots,a)}{n!}\right\|_{\rm rep}
   &\le C(J,n)A_n b^n L^{4-n},\notag\\
 \left\|\frac{(W{\bf1}_{\{D_P>T_0\}})(a,\ldots,a)}
                                      {n!}\right\|_{\rm rep}
   &\le C(J,n)A_n b^n L^{4-n}e^{-\mu T_0/2}.
 \label{fcc-tail-localization}
\end{align}
The target exponent is exactly \(a_0\). The extra
carrier length was paid after spatial compression
by \(JL\), rather than by lowering the old-scale
activity exponent.

Apply the second inequality to \(K_n^{\rm tail}\),
using \(A_n\le C(J)\varepsilon b^{-n}\) and
\(T_0\ge P/32\). Apply the first to the fixed-range
repair \(R_n\), retaining its factor \(e^{-c(J)P}\)
from \eqref{fcc-repaired-bounds}. Relative periodization
counts every occurrence, including any repeated
periodic slots. Together these give
\begin{equation}
 \left\|\sum_{n=2}^4
    \frac{(\operatorname{Per}_P R_n+K_n^{\rm tail})
                         (a,\ldots,a)}{n!}\right\|_{\rm rep}
 \le C_{\rm err}(J)\varepsilon
              (L^2+L+1)e^{-c(J)LN}.
 \label{fcc-error-estimate}
\end{equation}
All \(b\)-powers cancel in their correct degree.
An old-scale tail estimate that loses part of
\(a_0\) is not being invoked.

\subsection{High degrees and the finite quotient estimate}
\label{fcc-contraction-section}

Coefficientwise homogeneous phase projection has norm
at most one. Thus
\begin{equation}
 R_{\ge5}E=E-\sum_{n=0}^4P_nE,\qquad
 \|R_{\ge5}E\|\le6\|E\|.
 \label{fcc-high-projection}
\end{equation}
The supported fifth-order background theorem gives
\begin{equation}
 \|\mathcal T_{\rm high}R_{\ge5}E\|_{N;a_0,b}
        \le6C_{\rm high}\varepsilon/L,
 \qquad C_{\rm high}=C_0(8C_A)^5.
 \label{fcc-high-estimate}
\end{equation}
After the fixed classical reserves, \(C_{\rm high}\)
does not depend on \(M,L,J,N\). A fixed \(J\) satisfying
the earlier tail budgets can be used while increasing
the endpoint to \(L=R^q\): the core-compression
coefficient decreases, whereas the \(D(J)t\) and
projection budgets are unchanged. The field-scaling
disk has radius \(L/(8C_A)\); choose \(b\) afterward
so that
\begin{equation}
 b\le1,\quad b\le\rho/2,\quad
 bL/(8C_A)\le r_A,\quad L/(8C_A)\ge2,
 \label{fcc-radius}
\end{equation}
and impose the other fixed classical domain
restrictions.

The degree-zero sector disappears from the normalized
total. The degree-one total vanishes by semisimple
global color invariance and remains zero after
substitution, even if its individual coefficients
are nonzero. The degree-two through degree-four
polynomials are exactly
\eqref{fcc-polynomial-decomposition}. Thus the low
estimate, error estimate, and high estimate represent
the entire finite normalized substituted total,
not merely its first few kernels.

\begin{theorem}[Finite classical contraction]
\label{fcc-theorem}
Fix \(J\) and its stated classical tail and support
budgets. There are \(L_0(J)\), \(N_0(J)\), positive
\(c(J)\), and constants \(C_{\rm low}(J)\),
\(C_{\rm err}(J)\), \(C_{\rm ext}(J)\), independent
of \(L,b,M,N\), such that every allowed
\(L\ge L_0(J)\), every compatible \(N\ge N_0(J)\),
and every sufficiently small \(b\) chosen after \(L\)
satisfy
\begin{align}
 &\|\mathcal P_{ML,N}\mathcal T_{M,L,N}
                              \mathcal P_{M,LN}\|\notag\\
 &\quad\le C_{\rm ext}(J)^2
       \left[
       C_{\rm low}(J)L^{-\alpha}
       +\frac{6C_{\rm high}}L
       +C_{\rm err}(J)(L^2+L+1)e^{-c(J)LN}\right].
 \label{fcc-main-bound}
\end{align}
The input and output finite invariant quotients have
the same exponent \(a_0\) and radius \(b\).
Separate input and output projection constants can
replace the square if their conventions differ.
\end{theorem}

\begin{proof}
Take an almost-minimal representative of the extracted
class. Its norm can approach
\(\|\mathcal P_{M,LN}[E]\|\le C_{\rm ext}(J)\|[E]\|\).
Apply \eqref{fcc-canceled-estimate},
\eqref{fcc-error-estimate}, and
\eqref{fcc-high-estimate} to this representative.
They give an exactly equivalent connected output
family. The output projection costs the second
\(C_{\rm ext}(J)\). Take the quotient infimum and
send the representative error to zero.
\end{proof}

Renaming constants, the bound has the form
\begin{equation}
 \|\mathcal P_{\rm out}\mathcal T\mathcal P_{\rm in}\|
 \le A(J)L^{-\alpha}+B(J)/L
                       +E(L,J)e^{-c(J)N}.
 \label{fcc-renamed-bound}
\end{equation}
The sharper \eqref{fcc-main-bound} displays the actual
input period \(LN\). For any desired \(\theta>0\),
choose a finite allowed power \(L=R^q\) after the
fixed \(J\) to make the first two terms at most
\(\theta/2\); enlarge the compatible threshold
\(N_0\) until the error is at most \(\theta/2\) for
all \(N\ge N_0\); and choose \(b\) afterward.
This is a \(\theta\)-contraction uniform in preceding
depth and all these finite regulators. Shrinking
\(b\) does not change its constants.

By \eqref{fcc-telescope} the same bound applies to
\(\mathcal P_{\rm out}\mathcal T\), with the input
extraction implicit. That identity uses the full
actual finite reference, not merely an agreement
of flat Hessians.

\subsection{Scope}
\label{fcc-scope}

The estimate uses actual finite normalized totals
and periodic spatial algebra. It requires neither
an infinite-volume interaction limit nor evaluation
of an infinite total at a nonsummable periodic
field. Small finite lifted Ward defects and
wrapping tails are retained explicitly. The
resulting small representatives are exactly
equivalent on the full finite \(b\)-polydisc; the
original carrier assignment need not itself
obey the improved norm.

The order is fixed \(J\), then a finite endpoint
\(L=R^q\), then the small \(b\). This removes the
finite-regulator classical relocalization issue.
It does not identify \(q\) sequential quantum
cutoff integrals with one direct macro cutoff
integral. Application to an actual finite quantum
recursion still needs the exact finite scalar
descent, every intermediate analytic and Gaussian
budget, the running-coupling bounds and the
specified cutoff/history identities. Raw vacuum
normalizations remain separate. No beta-function
sign, indefinitely iterable estimate, large-field
control or continuum existence claim follows
from this classical result alone.
